\documentclass{article}
\usepackage[utf8]{inputenc}
\usepackage[T1]{fontenc}
\usepackage{amsmath}
\usepackage{amssymb}
\usepackage{array}
\usepackage{geometry}
\geometry{margin=1in}
\begin{document}
\section*{Genetic Algorithm (GA) for TSPTW-D}
\subsection*{Definition}
A Genetic Algorithm (GA) is a stochastic metaheuristic optimization technique inspired by natural selection. This
implementation addresses the \textbf{Traveling Salesman Problem with Time Windows and Dynamic costs (TSPTW-D)},
where the objective is to minimize the total duration of a tour while respecting city-specific availability
windows and accounting for time-dependent traffic perturbations.
\subsection*{Methodology}
The algorithm represents a potential tour as a chromosome and uses the following operators and models:
\begin{itemize}
    \item \textbf{Chromosome}: A permutation of indices representing the sequence of cities visited: $\sigma =
(\sigma_1, \sigma_2, \dots, \sigma_n)$.
    \item \textbf{Temporal Model}: Each city $i$ has an opening time $a_i$, a deadline $b_i$, and a service time
$s_i$.
    \item \textbf{Clock Propagation}: The arrival time $\tau$ at the next city $\sigma_{k+1}$ is calculated by:
    \[ \tau_{\sigma_{k+1}} = \max(\tau_{\sigma_k}, a_{\sigma_k}) + s_{\sigma_k} + c_{\sigma_k,
\sigma_{k+1}}(d_{\sigma_k}) \]
    where $d_{\sigma_k}$ is the departure time from city $\sigma_k$ and $c_{ij}(t)$ is the dynamic travel cost.
    \item \textbf{Dynamic Costs}: Travel time between $i$ and $j$ depends on the departure time $t$:
    \[ c_{ij}(t) = c_{base}[i][j] \times (1 + \delta_{ij}(t)) \]
    where $\delta_{ij}(t)$ represents a congestion factor (perturbation) active during specific intervals.
    \item \textbf{Fitness Function}: A soft-constraint approach penalizes late arrivals ($\tau_i > b_i$) using a
large coefficient $\lambda$:
    \[ F(\sigma) = \tau_0^{return} + \lambda \sum_{k=1}^{n} \max(0, \tau_{\sigma_k} - b_{\sigma_k}) \]
    \item \textbf{Operators}: Tournament selection, Ordered Crossover (OX), and 2-opt Mutation are used to evolve
the population while maintaining valid permutations.
\end{itemize}
\subsection*{Population Diversity}
To monitor the health of the genetic search, the population diversity is estimated by the standard deviation of
fitness scores:
\[ \sigma_{div} = \sqrt{\frac{1}{N}\sum_{i=1}^{N}\bigl(f(\pi_i) - \bar{f}\bigr)^2} \]
\subsection*{Code Example (2-opt Mutation)}
\begin{verbatim}
def mutation_2opt(chromosome, rate):
    chrom = chromosome[:]
    if random.random() < rate:
        i, j = sorted(random.sample(range(len(chrom)), 2))
        chrom[i:j+1] = reversed(chrom[i:j+1])
    return chrom
\end{verbatim}
\subsection*{Benchmark Results}
The table below summarizes the optimization results from the latest simulation runs (including penalties for
TSPTW-D cases):
\begin{center}
\begin{tabular}{|c|r|r|c|r|}
\hline
\textbf{Cities} & \textbf{Initial Score} & \textbf{Final Score} & \textbf{Gain (\%)} & \textbf{Time (s)} \\
\hline
10 (TSP) & 431 & 333 & -22.6\% & 0.99 \\
20 (TSPTW-D) & 74,391,170 & 29,229,845 & -60.7\% & 2.58 \\
100 (TSPTW-D) & 5,601 & 1,319 & -76.4\% & 56.54 \\
200 (TSP) & 10,203 & 1,238 & -87.9\% & 137.70 \\
500 (TSP) & 25,104 & 2,041 & -91.9\% & 1,794.99 \\
\hline
\end{tabular}
\end{center}
\textit{Note: The large scores for 20 cities reflect high penalties $\lambda=10^4$ applied to a synthetic
constrained dataset.}
\end{document}